% Self-contained research section for integration into the main report.
% Requires amsmath, amssymb, amsthm, hyperref, and theorem/lemma/proposition/
% corollary environments. No additional custom macros are used.
% New relative to the sixteen active ZIPs at ProveIt commit
% 4890fd00a5cacfc9881ae282120978f95c3ba392. No literature-priority claim.
% This is an obstruction theorem, NOT a below-frontier axiomatization.
\section{A stronger unary obstruction: full induction for quantifier-free formulas}
\label{sec:unary-amalgam}

The current unary obstruction leaves a specific question open: could
full quantifier-free induction with arbitrary numerical parameters repair
a language containing every unary primitive-recursive operation? The
answer is negative for the generous theory defined below. Its model
also has decidable equality and satisfies every true universal sentence,
not merely universal Horn sentences. Thus the failure persists without
the merging equality used in the earlier Kripke fan.

The construction is a classical union of elementary extensions, with
their standard elements identified. Unary operations can act within
one copy and can return to the shared standard part. They cannot join
nonstandard data from two different copies. A restriction argument
nevertheless validates \emph{every} quantifier-free induction matrix
with parameters. The gap is witnessed by one explicit existential
induction instance and by an automorphism obstruction to all ordinary
definable inverse interpretations.

\subsection{The exact source theory and interpretation convention}
\label{subsec:amalgam-source}

Let $L_1$ contain the numerical constant $0$ and a unary function symbol
for every fixed unary primitive-recursive function, including successor
$S$ and predecessor. It has no positive-arity relation symbols other
than equality, and no function of arity greater than one. Write
$\mathbb N_1$ for the standard structure. A numeral $\bar n$ denotes
the fixed term $S^n0$. A fixed composition of unary operations is an
ordinary term; there is no program variable or universal evaluator.

Quantifier-free formulas are built from equations and genuine falsity
by the usual propositional connectives. In the following axiom
families, \emph{true} means true in the classical standard structure
$\mathbb N_1$. This selects the axioms; the deductive logic of the
resulting theory is ordinary intuitionistic first-order logic with
actual equality and genuine falsity. In particular, we do not change
the meaning of implication or quantification when assessing an
interpretation.

Define $T_1$ to have the following nonlogical axioms.
\begin{enumerate}
\item \textbf{All true universal sentences.} Include every sentence
\[
  \forall x_1\cdots\forall x_k\,B(x_1,\ldots,x_k)
  \tag{U}\label{eq:amalgam-U}
\]
true in $\mathbb N_1$, where $B$ is quantifier-free and $k$ is any
fixed finite number. The matrix may use arbitrary propositional
connectives; it need not be Horn, positive, or an equation.
\item \textbf{All true one-input existence sentences.} Include every
true sentence
\[
 \forall x\,\exists y_1\cdots\exists y_m\,
                 C(x,y_1,\ldots,y_m),
 \tag{E1}\label{eq:amalgam-E1}
\]
where $C$ is quantifier-free, $m$ is fixed and finite, and the displayed
$x$ is its only universally quantified input. The free variables of
$C$ are among those displayed. Fixed standard parameters can be
written as numerals. Arbitrary additional universally quantified
parameters are deliberately not included in this family.
\item \textbf{Full quantifier-free induction.} For every
quantifier-free $A(x,\vec p)$, include the universal closure in
$\vec p$ of
\[
 \bigl(A(0,\vec p)\wedge
       \forall x\,(A(x,\vec p)\to A(Sx,\vec p))\bigr)
       \to\forall x\,A(x,\vec p).
 \tag{QFI}\label{eq:amalgam-QFI}
\]
The tuple $\vec p$ may have any fixed finite length. Equivalently,
use the usual contextual induction rule, keeping all parameters and
retained assumptions fixed and making its step variable fresh.
\end{enumerate}

This is a deliberately generous obstruction theory, not a proposed
effective presentation or a candidate proof checker. No recursive
enumerability claim is made for its all-true axiom families.
Its axioms already include successor injectivity and separation,
$0\ne S0$, every true quantifier-free fact with any finite number
of inputs, and the genuine two-variable equality-decidability sentence
\[
                 \forall x\forall y\,(x=y\vee\neg(x=y)).
 \tag{DecEq}\label{eq:amalgam-deceq}
\]
Numerical coverage follows from the instance
$\forall x\exists y\,(x=0\vee Sy=x)$ of \eqref{eq:amalgam-E1}.
It also follows from the true universal sentence
$\forall x\,(x=0\vee S\operatorname{pred}(x)=x)$.

An \emph{ordinary first-order interpretation} here may use a definable
subset of a finite Cartesian power, a definable equivalence relation
as its equality, and definable graphs of its functions. Logical
connectives retain their ordinary meanings, and quantifiers are
relativized in the usual way. Bi-interpretability additionally requires
provably definable isomorphisms to the two double interpretations.
The main theorem permits a fixed finite tuple of interpretation
parameters. A simple enlargement of the model, stated separately,
handles countably many parameters spread across an infinite signature.

\begin{theorem}[Unary amalgam obstruction]\label{thm:unary-amalgam}
The theory $T_1$ has a classical model with a nontrivial involutive
automorphism fixing any prescribed finite parameter tuple.
Consequently $T_1$ is not ordinarily bi-interpretable with iPRA,
even allowing finite-dimensional definable domains and definable
quotient equalities. The same conclusion holds for any theory in
this language whose axioms are true in the model constructed below,
including every subtheory of $T_1$.
\end{theorem}

\subsection{The model: copies sharing exactly the standard part}
\label{subsec:amalgam-model}

Choose a proper elementary extension
$M\succ\mathbb N_1$. One obtains it by adjoining to the elementary
diagram of $\mathbb N_1$ a constant different from every standard
numeral and applying compactness. Every finite set of the additional
requirements is realized by a sufficiently large standard numeral.
Since $L_1$ is countable, $M$ can also be chosen countable, although
countability is not needed below. Identify the named standard copy
of $\mathbb N_1$ with a substructure of $M$.

The phrase \emph{nonstandard} in this proof means externally that an
element belongs to $M\setminus\mathbb N$. We do not assert that
this set or the standard part is first-order definable in $M$.
The construction takes place in the ordinary classical metatheory.

Let $I$ be a countably infinite set of copy labels. Form the disjoint
union
\[
              D=\mathbb N\ \sqcup\bigl((M\setminus\mathbb N)\times I\bigr).
 \tag{D}\label{eq:amalgam-domain}
\]
Interpret $0$ by its standard value. For every unary symbol $f$, put
$f^D(n)=f^{\mathbb N_1}(n)$ at standard $n$, and put
\[
 f^D(a,i)=
 \begin{cases}
   f^M(a),&\text{if }f^M(a)\in\mathbb N,\\
   (f^M(a),i),&\text{if }f^M(a)\notin\mathbb N.
 \end{cases}
 \tag{F}\label{eq:amalgam-functions}
\]
Equality is actual identity in $D$ and falsity is false. These formulas
give total operations, including when a unary function takes a
nonstandard argument to a standard value.

For each $i\in I$, define the injection
$e_i:M\to D$ by $e_i(n)=n$ for standard $n$ and
$e_i(a)=(a,i)$ otherwise. Its image is denoted $M_i$.

\begin{lemma}[The copies are genuine substructures]
\label{lem:amalgam-copies}
Each $e_i$ is an isomorphism from $M$ onto the substructure $M_i$.
For $i\ne j$, one has $M_i\cap M_j=\mathbb N$, and
$D=\bigcup_{i\in I}M_i$. Inclusion of $M_i$ in $D$ preserves
and reflects the truth of every quantifier-free formula whose
parameters and arguments belong to $M_i$.
\end{lemma}
\begin{proof}
The operation law $f^D(e_i(a))=e_i(f^M(a))$ is exactly
\eqref{eq:amalgam-functions}, including its standard-output case.
The constant and injectivity conditions are immediate. The
intersection and union assertions follow from the disjoint union.
Injectivity and preservation of terms preserve and reflect equations;
in classical structures this extends through every propositional
connective and genuine falsity by induction on a quantifier-free
formula. No claim of elementarity of $M_i$ inside $D$ is made.
\end{proof}

\subsection{Why every true universal sentence survives}
\label{subsec:amalgam-universal}

Unary syntax is crucial here: a term contains at most one distinct
free variable. An equation can still compare terms in \emph{two}
independent variables, so preservation of all universal sentences is
stronger than checking each copy separately.

\begin{lemma}[Finite atomic diagrams can be copied into the naturals]
\label{lem:amalgam-diagram}
Fix a finite tuple $\vec a$ from $D$ and a finite set $\mathcal T$
of $L_1$ terms in its variables. There is a tuple $\vec b$ from
$\mathbb N$ such that for every $s,t\in\mathcal T$,
\[
          s^D(\vec a)=t^D(\vec a)
          \quad\Longleftrightarrow\quad
          s^{\mathbb N_1}(\vec b)=t^{\mathbb N_1}(\vec b).
 \tag{FD}\label{eq:amalgam-finite-diagram}
\]
Standard input entries can be left unchanged.
\end{lemma}
\begin{proof}
First fix every standard input entry to its original numeral. Let
$C\subseteq\mathbb N$ be the finite set of \emph{all} standard
values of terms from $\mathcal T$ at $\vec a$. Each remaining
input entry lies in a unique labelled copy outside the standard part.
Partition these inputs into their finitely many nonempty label groups.
Every term that is not already constant after the fixed substitutions
depends on inputs from just one group, because it has at most one
free variable. Assign that term to the group.

Process the groups in any order. Suppose groups before $i$ have
already been replaced by standard tuples, and let $V$ be the finite
set of term values produced by those previous choices. In the copy
with label $i$, forget its labels to obtain a tuple $\vec a_i$ in
$M$. Form a finite quantifier-free conjunction requiring the following:
\begin{enumerate}
\item every pair of terms assigned to this group has exactly the
same equality or inequality relation as at $\vec a_i$;
\item each assigned term with original standard value $c$ has value
exactly the numeral $\bar c$;
\item each assigned term with original nonstandard value differs from
every numeral for an element of $C\cup V$.
\end{enumerate}
Include equalities fixing any standard input variables used in the
group, or simply substitute their numerals before forming the conjunction.
All these requirements hold in $M$ at $\vec a_i$. In the last
family, this is because an original nonstandard value differs from
every specified standard number, regardless of how the preceding
standard choices were made.

The existential closure of this finite conjunction, with its finitely
many numeral parameters, is therefore true in $M$. Elementarity of
$\mathbb N_1\preccurlyeq M$ makes it true in $\mathbb N_1$.
Choose a standard tuple witnessing it, and proceed to the next group.
An empty collection of requirements causes no problem.

Within a group, all equations and inequalities are preserved by the
first condition. Terms with standard original values are fixed by the
second. A cross-group equality in $D$ is possible only when both
outputs are the same shared standard value; these equalities are
preserved. Every other cross-group comparison is an inequality. If
one output was standard and the other nonstandard, the latter was made
to avoid $C$. If both were nonstandard, the later group was made to
avoid the earlier group's actual selected output in $V$. Thus all
cross-group inequalities are preserved as well, proving
\eqref{eq:amalgam-finite-diagram}.
\end{proof}

\begin{proposition}[Universal theory preservation]
\label{prop:amalgam-universal}
The structure $D$ satisfies every axiom of family
\eqref{eq:amalgam-U}.
\end{proposition}
\begin{proof}
Suppose $\forall\vec x B(\vec x)$ is true in $\mathbb N_1$
but false in $D$, and choose a tuple witnessing its failure. Apply
Lemma~\ref{lem:amalgam-diagram} to the finitely many terms occurring
in its atomic subformulas. The selected standard tuple gives every
atom the same truth value. Induction on the quantifier-free matrix
$B$, including implication and negation, gives it the same classical
truth value, contradicting the original universal sentence.
\end{proof}

\subsection{Existence and full quantifier-free induction}
\label{subsec:amalgam-induction}

\begin{proposition}[All one-input existence axioms]
\label{prop:amalgam-one-input}
The structure $D$ satisfies every instance of
\eqref{eq:amalgam-E1}.
\end{proposition}
\begin{proof}
Given $a\in D$, choose a copy $M_i$ containing it. By elementarity,
$M_i$ satisfies the displayed one-input sentence and supplies a finite
tuple of witnesses inside $M_i$. Its matrix remains true in $D$
by Lemma~\ref{lem:amalgam-copies}. When $a$ is standard, any copy
may be used. This proof specifically uses the fact that there is only
one independently chosen input; no single copy need contain two
independent nonstandard inputs.
\end{proof}

\begin{lemma}[Restriction of a quantifier-free matrix]
\label{lem:amalgam-restrict}
Let $A(x,\vec p)$ be quantifier-free and fix arbitrary parameter
values $\vec p\in D$. For every label $i$, there is a
quantifier-free formula $A_i(x,\vec q_i)$ with parameters in
$M_i$ such that
\[
 D\models A(a,\vec p)
 \quad\Longleftrightarrow\quad
 M_i\models A_i(a,\vec q_i)
 \qquad(a\in M_i).
 \tag{Res}\label{eq:amalgam-restriction}
\]
No new function symbols or quantified predicates are required.
\end{lemma}
\begin{proof}
After evaluating parameter-only terms, every equation is of one of
three types. It compares two unary terms of $x$; it compares a unary
term of $x$ with a fixed element $b\in D$; or it compares two
fixed elements. In the first case retain the equation. In the
second, retain it with $b$ as a parameter when $b\in M_i$;
when $b\notin M_i$, replace it by $0=S0$. This replacement is
valid because every unary term of an element of $M_i$ remains in
$M_i$, and $0=S0$ is false in $M_i$. In the third case replace
the equation by $0=0$ or $0=S0$ according to its fixed truth value
in $D$. Use the same replacements when the two sides appear in
the reverse order. Only finitely many new parameter values are used.

Lemma~\ref{lem:amalgam-copies} verifies the retained equations;
the stated cases verify the replacements. Keep genuine falsity and
all propositional connectives unchanged. Induction on the formula
then proves \eqref{eq:amalgam-restriction}. The replacements depend
on the selected model parameters, which is permitted in a semantic
verification of induction; no uniform syntactic decision procedure is
being claimed.
\end{proof}

\begin{proposition}[Full parameterized quantifier-free induction]
\label{prop:amalgam-qfi}
The structure $D$ satisfies every instance of
\eqref{eq:amalgam-QFI}, including its contextual rule form.
\end{proposition}
\begin{proof}
Fix $A$ and $\vec p$ and assume in $D$ that
$A(0,\vec p)$ and
$\forall x(A(x,\vec p)\to A(Sx,\vec p))$ hold.
For each $i$, use Lemma~\ref{lem:amalgam-restrict}. Since
$0\in M_i$ and $M_i$ is closed under successor, the two global
premises restrict to the base and successor-step premises of
$A_i(x,\vec q_i)$ in $M_i$.

Every quantifier-free induction instance is true in $\mathbb N_1$:
external ordinary induction proves it for each fixed formula and
each fixed tuple of natural parameters. Its universal closure in
those parameters is therefore true there, and by elementarity it
holds in $M_i$. Applying it gives $A_i(a,\vec q_i)$ for every
$a\in M_i$, hence $A(a,\vec p)$ throughout that copy.
As the copies cover $D$, the conclusion holds throughout $D$.
This proves the universally closed induction axiom for arbitrary
$\vec p$.

For the contextual rule, fix a valuation satisfying the retained
context. If the two premises hold under that valuation, the same
argument applies. Freshness of the induction variable ensures that
the retained context and parameters do not change when the variable
ranges over a copy. Alternatively, the contextual rule follows from
the just-verified induction axiom by ordinary intuitionistic logic.
\end{proof}

Together these propositions make $D$ a classical model of $T_1$.
A classical structure regarded as a one-world intuitionistic Kripke
model has the same truth values for all formulas: at the only world,
the clauses for implication and the quantifiers reduce to their
classical clauses. Thus this model is a sound countermodel for the
specified intuitionistic theory, not only for a classical variant of
it. Its equality is actual identity and is decidable. No transition
maps, future witnesses, or identification of previously distinct
elements are involved.

\subsection{Why definable quotients cannot restore a bi-interpretation}
\label{subsec:amalgam-biinterpretation}

Every permutation $\rho$ of $I$ induces a permutation $\sigma_\rho$
of $D$: it fixes each standard element and sends $(a,i)$ to
$(a,\rho(i))$. Formula~\eqref{eq:amalgam-functions} shows directly
that it commutes with every unary function, in both the standard-output
and nonstandard-output cases. It is therefore an automorphism.

\begin{lemma}[Parameter-fixing involutions]
\label{lem:amalgam-involution}
For every finite tuple of elements of $D$, there is a nonidentity
automorphism of order two fixing that tuple pointwise.
\end{lemma}
\begin{proof}
The nonstandard entries of the tuple use only finitely many labels.
Select distinct unused labels $i,j$ and transpose them. The resulting
automorphism fixes the tuple and squares to the identity. If
$a\in M\setminus\mathbb N$, it moves $(a,i)$ to $(a,j)$, so
it is nonidentity. No cardinal counting of complete types is needed.
\end{proof}

The one fact about the target used below is that iPRA defines a linear
order on its whole numerical domain and proves its order laws. For
example, let $x\le y$ abbreviate $x\mathbin{\dot-}y=0$,
where $\mathbin{\dot-}$ is the usual fixed primitive-recursive truncated
subtraction operation. Its reflexivity, antisymmetry, transitivity,
and totality are elementary iPRA theorems proved using defining
equations and quantifier-free induction. This is the same target
property used in Section~2, \emph{The common rigidity argument}, of the
current \texttt{ipra\_unary\_barriers} report.

\begin{lemma}[Order-preserving involutions are trivial]
\label{lem:amalgam-order}
An order-preserving permutation $\tau$ of a linear order with
$\tau^2=\mathrm{id}$ is the identity.
\end{lemma}
\begin{proof}
If $a<\tau(a)$, applying $\tau$ gives
$\tau(a)<\tau^2(a)=a$, a contradiction. If $\tau(a)<a$,
applying $\tau$ gives the opposite contradiction. Linearity leaves
only $\tau(a)=a$. The proof applies to any linear order, finite
or infinite, standard or nonstandard.
\end{proof}

\begin{proof}[Proof of Theorem~\ref{thm:unary-amalgam}]
The preceding propositions and Lemma~\ref{lem:amalgam-involution}
establish the asserted model and automorphisms. Suppose now that
ordinary bi-interpretations with iPRA were given. Interpret iPRA in
$D$, obtaining a structure $N^*$ with underlying set
$X/E$, where $X\subseteq D^d$ and $E$ is the interpreted equality.
The translated iPRA axioms and the well-definedness obligations for
this interpretation hold in $D$, because they are provable in the
source theory.

Collect the finite source parameter tuple used by the interpretation,
the reverse interpretation, and the source-definable composite
isomorphism. A parameter in $N^*$ is an $E$-class; choose a
representing tuple in $X$ and include its finitely many coordinates
among the source parameters. The same procedure handles any further
quotient parameters used for the composite. When parameters are
specified by a provable existential parameter condition, choose its
witnesses in $D$ first and include them as well. This is still a
finite parameter set under the stated interpretation convention.

Choose an involution $\sigma$ fixing all those elements.
Applying it componentwise preserves $X$, $E$, and every defining
relation and function graph. It induces the well-defined map
\[
           \sigma^*([\vec a]_E)=[\sigma(\vec a)]_E.
 \tag{Ind}\label{eq:amalgam-induced}
\]
The map is an automorphism of $N^*$, and
$(\sigma^*)^2=\mathrm{id}$ because the construction respects
composition. The interpreted order laws make the interpreted order
an actual linear order on $X/E$ in this classical model.
Lemma~\ref{lem:amalgam-order} therefore gives
$\sigma^*=\mathrm{id}$ on $N^*$.

Apply the reverse interpretation to $N^*$. Since $\sigma^*$
is already the identity and fixes its parameters, its induced map
$\sigma^{**}$ on the double interpreted domain is the identity.
This remains so when that domain is a definable subset of a power
or a quotient of such a subset: componentwise identity induces
identity on every class.

By bi-interpretability, the source theory proves that a fixed
definable graph is an isomorphism $j:D\to D^{**}$. Its
interpretation in $D$ is therefore a total injective function to
the double interpreted equality quotient. Invariance of that graph
under the parameter-fixing automorphism gives
\[
                   j(\sigma a)=\sigma^{**}(j(a))=j(a).
 \tag{Nat}\label{eq:amalgam-natural}
\]
Injectivity forces $\sigma a=a$ for every $a\in D$, contrary
to the chosen moved element. This proves the obstruction. The same
argument uses only that $D$ is a model of the proposed source theory,
so it applies to the additional theories stated in the theorem.
\end{proof}

\paragraph{Optional parameter extension.}
The main theorem needs only infinitely many labels and finite global
parameter tuples. If one permits countably many parameters spread
across the infinitely many symbol definitions, choose $I$ uncountable
instead. Their representatives occupy at most countably many labels,
so two unused labels remain. More generally, for any prescribed
infinite cardinal bound $\lambda$ on all such parameters, take
$|I|>\lambda$; finite-dimensional representatives do not increase
that bound. The preceding proofs of the axioms do not depend on
$I$ being countable. These observations handle the broader parameter
conventions without changing the essential finite-parameter theorem.

\subsection{A concrete existential-induction gap}
\label{subsec:amalgam-pairing}

The failure can be exhibited by fixed familiar unary functions.
Choose a primitive-recursive bijective pairing on the naturals, for
example
\[
          \operatorname{pair}(x,y)=2^x(2y+1)-1,
\]
and let $\pi_0,\pi_1$ be its unary primitive-recursive inverse
projections. Their bounded search definitions make them fixed unary
primitive-recursive functions and hence members of $L_1$.
The displayed binary pairing is used here in the metatheory to define
those unary functions; it is \emph{not} inserted as a source symbol.

\begin{proposition}[Failure of pairing totality]
\label{prop:amalgam-no-pair}
If $a\in M_i\setminus\mathbb N$ and
$b\in M_j\setminus\mathbb N$ with $i\ne j$, then
\[
 D\not\models\exists z\,
              (\pi_0(z)=a\wedge\pi_1(z)=b).
 \tag{NoPair}\label{eq:amalgam-no-pair}
\]
In particular, the true standard sentence
\[
 \forall x\forall y\exists z\,
              (\pi_0(z)=x\wedge\pi_1(z)=y)
 \tag{PairTot}\label{eq:amalgam-pairtot}
\]
is not derivable in $T_1$.
\end{proposition}
\begin{proof}
If $z$ is standard, both unary images are standard. If $z$ has
label $k$, every nonstandard unary image has that same label $k$,
by \eqref{eq:amalgam-functions}. Its two nonstandard images cannot
simultaneously have the distinct labels $i,j$. The failure in a
model of $T_1$ proves nonderivability by intuitionistic soundness.
\end{proof}

Now define the following two \emph{unary} primitive-recursive
functions, which are already in $L_1$:
\[
 c(p)=\operatorname{pair}(0,p),\qquad
 d(z)=\operatorname{pair}(S\pi_0(z),\pi_1(z)).
\]
The four true universal equations
\[
 \begin{aligned}
 \pi_0(c(p))&=0,&\pi_1(c(p))&=p,\\
 \pi_0(d(z))&=S\pi_0(z),&\pi_1(d(z))&=\pi_1(z)
 \end{aligned}
 \tag{PC}\label{eq:amalgam-pair-compute}
\]
are axioms of $T_1$. Define the genuinely existential matrix
\[
 A(n,p):=\exists z\,
             (\pi_0(z)=n\wedge\pi_1(z)=p).
 \tag{A}\label{eq:amalgam-existential-matrix}
\]

\begin{corollary}[One explicit induction instance exceeds the source]
\label{cor:amalgam-induction-gap}
The theory $T_1$ proves both $\forall p\,A(0,p)$ and
$\forall n\forall p(A(n,p)\to A(Sn,p))$, but does not prove
$\forall n\forall p\,A(n,p)$. Thus it does not admit induction
on the single $\Sigma_1$ matrix
\eqref{eq:amalgam-existential-matrix}, despite admitting all
quantifier-free induction with arbitrary parameters.
\end{corollary}
\begin{proof}
For the base choose the witness $c(p)$ and use the first two
equations of \eqref{eq:amalgam-pair-compute}. For the step,
eliminate a witness $z$ to $A(n,p)$ and choose the new witness
$d(z)$. The last two equations, equality transport, and conjunction
introduction prove its two required equations. Existential
introduction and elimination give the step in ordinary intuitionistic
logic, with the witness variable fresh for the conclusion.
The asserted universal conclusion is exactly
\eqref{eq:amalgam-pairtot}, which fails by
Proposition~\ref{prop:amalgam-no-pair}.
\end{proof}

This isolates an induction obstacle in reconstructions based on
function graphs. An equation mentioning graph-defined operations
can translate to an existential matrix in the old language; its
induction instance therefore needs a separate admissibility argument.
Here the pairing relation exhibits a concrete existential matrix with
a proved base and step but without its universal conclusion. The
example concerns the failure of two-input pairing totality itself;
no total pairing graph has been assumed. We do not claim that adding
just this one existential
induction instance would by itself make $T_1$ bi-interpretable with
iPRA; the positive reconstruction would still require its own proof.

\subsection{Exact comparison with the current corpus and scope}
\label{subsec:amalgam-comparison}

The comparison is with
\texttt{ipra\_unary\_barriers\_2026-11-09.zip}, present in
\texttt{docs/drafts} at commit
\path{4890fd00a5cacfc9881ae282120978f95c3ba392}.
Its Section~3, \emph{A pure-unary obstruction beyond monotonicity},
is the source file \path{sections/03_unary_fan.tex}. That section
uses a homomorphic Kripke fan whose root consists of vectors with
homogeneous complete one-types. Its source theory has all true
universal Horn sentences, all true parameter-free unary
quantifier-free sentences, coverage, equation induction with
parameters, and parameter-free quantifier-free induction.

The present theorem strengthens that source in three precise ways.
First, family \eqref{eq:amalgam-U} allows every true universal
sentence with any finite number of inputs and any quantifier-free
matrix. It contains the earlier Horn and one-variable families.
In particular, it includes \eqref{eq:amalgam-deceq}, whose failure
was exhibited at the earlier fan's root. Second,
\eqref{eq:amalgam-E1} includes every true quantifier-free existence
law with one input, of which coverage is only one example.
Third, \eqref{eq:amalgam-QFI} allows all quantifier-free matrices
with arbitrary finite parameter tuples, extending both induction
families of the earlier theorem. These are changes to the
mathematical scope, not reductions in source-code length.

The ordinary-interpretation convention and the target's linear-order
rigidity argument are inherited in method from Section~2, \emph{Target,
interpretations, and semantic conventions}, particularly its
subsection \emph{The common rigidity argument}. The new part is the
amalgam model and the proofs that it validates the larger axiom
families. Section~4 of the existing report concerns a distinct
zero-initialized recursion language; its monotonicity and fan results
are neither assumed nor replaced here. No claim of historical
priority beyond this inspected corpus is made.

The obstruction has deliberately stated limits. A binary operation
that combines independently chosen inputs, the pairing-totality
sentence \eqref{eq:amalgam-pairtot}, or another principle false in
$D$ can invalidate it. An interpretation that changes the meanings
of the logical connectives also lies outside the ordinary
first-order argument. The theorem does not prohibit all language or
domain changes allowed by the task, and it says nothing about
equiconsistency alone. It supplies a stronger test that a proposed
unary reduction must pass. It does not supply the requested
axiomatization or a complete checker below the current observed
frontier.

In particular, the successful closed-unary-update theories in the
current corpus are not counterexamples to this result. Their operators
$I_h(n,a)$ or $K_h(n,a)$ have two independent numerical arguments,
even though the fixed update description $h(u)$ is unary. The present
theorem concerns an entirely unary first-order numerical signature.
A binder that captures a numerical parameter can likewise conceal
additional first-order arguments and must not be classified as unary
merely because it displays one local update variable.

